\originalchapter{Boltzmann碰撞与H定理}
\label{ch:boltzmann}
\source{8.333 L8--L10。}
\originalsection{两体碰撞的增益与损失}
令$f(r,p,t)$归一为粒子数。稀薄经典气的局域二体碰撞项为
\begin{equation}
C[f](p_1)=\int\dd p_2\dd\Omega\,|v_1-v_2|\frac{\dd\sigma}{\dd\Omega}\,[f(p_1')f(p_2')-f(p_1)f(p_2)].
\end{equation}
$\sigma$是两体微分截面，$p_1',p_2'$与$p_1,p_2$由动量、能量守恒及散射角联系。第一项为碰撞入当前动量的增益，第二项为离开的损失。微观可逆性使正反散射的权重及相空间测度相同。Boltzmann方程为$\partial_tf+v\cdot\nabla_rf+F\cdot\nabla_pf=C[f]$。

量子气在末态还需乘$1\pm f$（占据数按单态归一），分别表示玻色增强和Pauli阻塞。不能直接把经典$\log f$的H函数拿来声称量子H定理；相应量子熵要包含$(1\pm f)\log(1\pm f)$。
\originalsection{碰撞不变量}
对任意函数$\psi(p)$，交换粒子1、2并将正反过程成对，
\begin{equation}
\int\psi_1C[f]\dd p_1=\frac14\int\dd\mathcal C\,(f_1'f_2'-f_1f_2)(\psi_1+\psi_2-\psi_1'-\psi_2'),
\end{equation}
$\dd\mathcal C$含非负散射测度。对$\psi=1,p_i,p^2/(2m)$，最后括号为0，碰撞分别保持粒子数、总动量和动能。这给流体守恒方程的微观来源。一般外场在输运左侧改变动量或能量，不是碰撞守恒失败。
\originalsection{H定理与平衡态}
定义$H[f]=\int f\log f\dd r\dd p$。无净边界熵流时输运项对$H$积分为0。利用粒子数守恒去掉$+1$导数项，碰撞贡献化为
\begin{equation}
\frac{\dd H}{\dd t}=-\frac14\int\dd r\dd\mathcal C\,(f_1'f_2'-f_1f_2)\log\frac{f_1'f_2'}{f_1f_2}\le0.
\end{equation}
因为对正$x,y$有$(x-y)\log(x/y)\ge0$，每个碰撞通道非正。物理粗粒度熵为$-\kb H$加计数常数，所以增加。该不等式是在Boltzmann闭合之下的证明，不是对任意精确$\rho_N$的声明。

平衡要求每条可达碰撞$f_1f_2=f_1'f_2'$，故$\log f$为碰撞不变量。对能连接整个动量空间的普通弹性散射，$\log f=a+\vct b\cdot p-cp^2$，完成平方得到
\begin{equation}
f_{\mathrm{le}}(r,p)=n(r)[2\pi m\kb T(r)]^{-d/2}\exp[-|p-mu(r)|^2/(2m\kb T(r))].
\end{equation}
这是局域Maxwell分布，碰撞项局域为0，但若$n,T,u$随位置变化，输运项一般不为0；全局平衡还需满足外势与边界条件。若碰撞有额外守恒量或不能混合所有动量，其零空间可能更大。
\originalsection{可逆性问题的具体边界}
Loschmidt反转将所有粒子速度倒转，同时将出射碰撞相关转成未来的入射相关，此特殊准备不满足分子混沌。Poincaré复现适于有限、封闭、有限相空间测度体系，指出无穷时间永久单调的精确熵增不能普遍成立。两者要求辨别热力学极限、可观测时间与统计初态，不能用“粒子多”一句话把数学反例抹去。

Boltzmann方程因忽略细相关而获得有效不可逆性。当观测尺度远大于碰撞时间，闭合可准确预测输运，即使全微观态的相关信息仍在增长。关于该近似的严格适用时间存在更深的数学问题，本稿未证明任意长时间、任意密度的Boltzmann极限。
